\documentclass{article}
\usepackage{geometry}
\geometry{a4paper, margin=1in}
\usepackage{amsmath}
\usepackage{amssymb}
\usepackage{graphicx}
\usepackage{kotex}
\usepackage{hyperref}

\title{ZUZU와 Claude Code: 1인 · 양국(兩國) 소프트웨어 회사}
\author{문현재 (Hyun Jae Moon) \\ \small 소프트웨어 엔지니어 \\ \small \href{mailto:hyunjae@moonai.co.kr}{hyunjae@moonai.co.kr}}
\date{2026년 7월 15일}

\begin{document}

\maketitle

\begin{abstract}
이 논문은 저자의 두 시기, 즉 구글에서의 3년간(2022--2025) 소프트웨어 엔지니어링과, 2025년부터 진행한 문AI(Moonai)의 1인 구축을 바탕으로 한 경험 보고서다. 문AI는 창업자가 이끄는 AI 납품 회사로, 2025년에 개발되어 2026년에 공식 출범했으며, 두 관할권(대한민국과 미국)에 걸쳐 단 한 사람이 운영한다. 두 가지 관찰이 이 보고서를 관통한다. 첫째, 1인 창업자는 포트폴리오의 모든 제품에 대해 완전한 애플리케이션 개발(소비자 수준의 앱이 요구하는 지속적인 프론트엔드 · 모바일 · 제품 디자인 반복)을 현실적으로 제공할 수 없다. 이 역량의 공백은 실재하며 사라지지 않았다. 둘째, 더 놀랍게도 두 개의 인프라가 그럼에도 Google Cloud Platform(GCP)과 Amazon Web Services(AWS) 양쪽에서 프로덕션 백엔드 시스템 소프트웨어를 납품하기에 \emph{충분}했다. 백엔드 납품의 엔지니어링 비용을 붕괴시킨 에이전틱 코딩 도구(Claude Code)와, 한국 법인 운영의 행정 비용을 붕괴시킨 법인 운영 플랫폼(ZUZU, \href{https://zuzu.network}{zuzu.network})이다. 세 번째 조건, 곧 창업자가 대한민국 국민이자 동시에 미국 영주권자(green card holder)라는 이중의 발판이 두 나라에서 동시에 사업을 운영하는 것을 법적 · 실무적으로 가능하게 했다. 우리는 이 세 요소가 어떻게 결합하는지, 무엇을 가능하게 했는지, 무엇을 끝내 해결하지 못했는지, 그리고 이것이 에이전틱 AI 시대의 소프트웨어 회사가 갖춰야 할 최소 형태에 대해 무엇을 시사하는지를 서술한다.
\end{abstract}

\section{서론}

2010년대의 통념은 소프트웨어 회사가 팀 스포츠라는 것이었다. 백엔드 엔지니어, 프론트엔드 엔지니어, 모바일 개발자, 디자이너, 운영 담당자, 그리고 설립 · 신고 · 준법의 법적 · 행정적 기계장치를 다룰 사람. 이 기능들 가운데 하나라도 빠진 창업자는 그것을 채용하거나, 아니면 실패할 것으로 기대되었다.

2022년부터 나는 그 구조 가운데 실제로 무엇이 회사를 지탱하는지에 관한 실험을, 대체로 자의가 아닌 상황 속에서 수행해 왔다. 2025년까지 구글에서의 3년간 소프트웨어 엔지니어링(그 상당 부분은 네트워크 시뮬레이션 소프트웨어의 릴리스 엔지니어링과 CI/CD였다), 그리고 2025년부터의 문AI 1인 구축이 그것이다. 문AI(한국어로는 \emph{문AI}, `문'은 한글로 표기)는 한국 중소기업(SMB)을 겨냥한 AI 납품 회사이며, 구성원은 정확히 한 명, 나뿐이다. 2025년은 개발의 해였고, 2026년은 공식 출범의 해, 곧 전체 과정(미국 LLC 설립, 한국 법인 설립, 공급기업 등록, 그리고 첫 유료 엔게이지먼트)이 하나로 모이는 시점이다. 회사를 만들면서 나는 구글에서의 시기가 벼려 준 상근 엔지니어링 관점을 바탕으로 제품 포트폴리오를 세우려 했다. AI 기반 산업용 카탈로그(yushin-melamine.com), 소셜 게임 트래킹 플랫폼(playshelfapp.com), 자율 파이썬 테스팅 에이전트(pythontestingagent.com), 소비자향 AI 경험들, 그리고 납품 사업 그 자체다.

솔직한 결과는 비대칭적이다. 나는 여러 문AI 제품에 대해 제품 규모의 앱 개발을 \emph{제공하지 못했다}. 소비자 애플리케이션 개발의 다듬기 반복, 곧 디자인 반복, 플랫폼별 모바일 작업, 데모를 사람들이 실제로 머무는 제품으로 갈라놓는 매일의 UI 정제 노동은 한 사람이 포트폴리오 전반에 걸쳐 지속할 수 있는 범위를 넘어선다. 그러나 나는 GCP와 AWS 양쪽에서, 배포되고 모니터링되며 다중 리전으로 실제 고객을 서빙하는 프로덕션급 백엔드 시스템 소프트웨어를 \emph{제공할 수 있었다}. 두 결과의 차이는 재능이나 노력이 아니었다. 그것은 일의 한쪽에는 특정 인프라가 존재했고 다른 쪽에는 존재하지 않았다는 사실이었다.

이 논문은 그 두 조각(엔지니어링 납품의 Claude Code와 법인 운영의 ZUZU)을 명명하고, 그것들이 왜 충분했는지 살피며, 내가 대한민국 국민이자 미국 영주권자라는 이중 지위가 가능케 한 국경 횡단 구조 안에 그것들을 자리매김하고, 다음의 교훈을 일반화한다. 최소 형태의 소프트웨어 회사는 이제 그 기능들 가운데 어느 것이 충분한 인프라 속으로 흡수되었고 어느 것이 아직 아닌지에 의해 정의된다.

\section{배경: 한 명의 창업자, 두 나라, 하나의 포트폴리오}

\subsection{국경을 넘는 위치}

나는 대한민국 국민이며 미국의 합법적 영주권자(green card holder)다. 내 경력의 대부분에서 행정적 부담(두 세제, 두 벌의 신고, 미국 쪽의 거주 의무)처럼 느껴졌던 이 조합은, 알고 보니 뒤따르는 모든 것의 구조적 전제 조건이었다.

한국 국적은 나를 한국 비즈니스 생태계의 일급 참여자로 만든다. 나는 한국 법인을 설립하고, 정부 보조 디지털화 사업의 공급기업으로 등록하며, 외국 법인에게는 사실상 닫혀 있는 조달 · 보조 체계 안에서 한국 중소기업에 판매할 수 있다. 미국 영주권은 나를 미국 생태계의 일급 참여자로 만든다. 나는 캘리포니아에 거주하고, 1일차 청구 수단으로 단일 구성원 LLC를 세우며, 미국 은행 및 결제 레일을 열고, 비자 리스크 없이 미국 고객과 계약할 수 있다.

어느 한 나라의 위치만으로는 문AI의 모델을 지탱할 수 없다. 시장과 납품 플레이북은 한국의 것이고, 창업자의 거주지 · 초기 현금흐름 · 클라우드 청구 관계는 미국의 것이다. 양쪽 모두에 법적으로 토착적이라는 사실이, 1인 회사가 관할권 측면에서 마치 국제적 자회사 구조를 갖춘 회사처럼 행동하게 한다.

\subsection{포트폴리오와 역량의 경계}

문AI의 역량 증명 포트폴리오는 여러 프로덕션 시스템에 걸쳐 있다. 플래그십은 한국 멜라민 제조사를 위한 AI 기반 이중언어 제품 카탈로그로, GCP에서 돈다. \texttt{us-central1}과 \texttt{asia-northeast3}의 Cloud Run 서비스, 검색과 채팅을 위한 Vertex AI Gemini, 카탈로그 데이터를 위한 Firestore, 자산을 위한 Google Cloud Storage, 파일 메타데이터를 위한 BigQuery. 회사 자체의 웹 프레즌스는 AWS에서 돈다. S3와 Lambda 함수 URL 앞에 놓인 CloudFront, 그리고 지속성을 위한 관리형 Postgres로, 프리 티어 안에 들어앉도록 선택한 아키텍처다. 납품 플레이북 자체가 명시적으로 듀얼 클라우드다. 동일한 모듈 라이브러리가 GCP에서는 Cloud Run + Vertex AI + Firestore + GCS로, AWS에서는 App Runner/Fargate + Bedrock + DynamoDB + S3로 매핑된다.

포트폴리오가 \emph{담지 못한} 것은 내가 인력으로 충당할 수 없었던 것, 곧 지속적인 애플리케이션 개발이다. 소셜 게임 트래킹 제품 Playshelf에는 실제 백엔드가 있다(FastAPI, 마이그레이션을 갖춘 PostgreSQL, 레이트 리밋을 인지하는 백그라운드 워커). 그러나 그 소비자향 애플리케이션 표면은 자금을 갖춘 소비자 스타트업이 요구할 속도의 일부로만 나아갔고, 지금도 초대제로 남아 있다. 몇몇 제품 구상은 프론트엔드를 아예 갖추지 못했다. 이 패턴은 우연이 아니라 하나의 법칙이라 부를 만큼 반복되었다. \textbf{1인 창업자로서 백엔드 시스템 소프트웨어는 납품 가능했고, 제품 규모의 앱 개발은 불가능했다.}

흥미로운 질문은 왜 그 선이 정확히 거기에 그어졌는가다.

\section{Claude Code: 에이전틱 코딩, 사라진 백엔드 팀}

\subsection{무엇이 달라졌나}

첫 번째 답은 이 시기에 성숙해 가는 동안 내가 채택한 에이전틱 코딩 도구 Claude Code다. 여기서 AI 코드 어시스턴트와 에이전틱 코딩 도구의 구분이 중요하다. 어시스턴트는 타이핑하는 사람을 가속하고, 에이전트는 사람의 키 입력이 아니라 사람의 지시 아래 다단계 엔지니어링 작업을 실행한다. 코드베이스를 읽고, 여러 파일에 걸쳐 코드를 쓰고, 명령을 실행하고, 오류를 읽고, 반복한다.

백엔드 시스템 작업은 그 일의 본질적 성질 때문에 이 변화가 가장 극적으로 나타나는 영역이었다.

\begin{itemize}
    \item \textbf{백엔드 정확성은 기계 검증이 가능하다.} Cloud Run 서비스는 요청을 서빙하거나 못 하거나 둘 중 하나다. IAM 정책은 호출을 허용하거나 403을 반환한다. 마이그레이션은 적용되거나 실패한다. 에이전트가 자율적으로 반복하기 위해 필요한 피드백을 시스템 자신이 즉각적이고 명확하게 생성한다. 배포 스크립트, 인프라 인벤토리, API 계약이 그 루프를 닫는다.
    \item \textbf{클라우드 플랫폼은 밑바닥까지 텍스트다.} GCP와 AWS는 CLI, 선언적 구성, 잘 문서화된 SDK로 운영되며, 이는 정확히 언어 모델이 가장 강한 매체다. 두 공급자의 서비스는 서로 준-동형이다(Cloud Run $\leftrightarrow$ App Runner, Firestore $\leftrightarrow$ DynamoDB, Vertex AI $\leftrightarrow$ Bedrock, GCS $\leftrightarrow$ S3). 그래서 한쪽에서 시스템을 구축한 에이전트를 다른 쪽으로 매핑하도록 지시할 수 있다. 문AI의 ``하나의 플레이북, 두 클라우드'' 원칙이 1인에게 경제적으로 성립하는 것은 오직 그 매핑 작업이 수작업이 아니라 에이전틱하기 때문이다.
    \item \textbf{창업자의 판단이 여전히 희소한 입력이며, 백엔드 판단은 잘 압축된다.} 내 역할은 코드를 쓰는 데서 아키텍처를 명세하고, diff를 검토하며, 보안과 비용 자세(시크릿 관리, 인증 패턴, 리전 전략, 프리 티어 규율)를 소유하는 쪽으로 옮겨 갔다. 경험 있는 한 사람의 판단이 에이전트로 증폭되어 작은 포트폴리오 전체의 백엔드 표면을 감당한다.
\end{itemize}

구체적 결과는 이렇다. 다중 리전 배포 템플릿, 검색 증강 생성(RAG) 파이프라인, 시크릿 매니저로 뒷받침되는 JWT 인증, 이미지 파이프라인, 인제스션 크롤러, 서버리스 웹 인프라가 두 클라우드에 걸쳐 한 사람에 의해 구축 · 배포 · 운영되었다. 몇 년 전이라면 이것은 3\,\textasciitilde 5인 엔지니어링 팀의 산출물이었다. 이것이 Claude Code가 \emph{충분}했다는 말의 정확한 의미다. 모든 줄을 그것이 썼다는 뜻이 아니라, 한 명의 창업자 더하기 한 개의 에이전틱 도구라는 조합이 문AI의 고객이 요구한 백엔드 납품 기준선 전체를 통과시켰다는 뜻이다.

\subsection{같은 도구가 앱 개발 격차를 메우지 못한 이유}

애플리케이션 쪽에서도 같은 레버리지를 기대하기 쉽지만, 5년치 기록은 그렇지 않다고 말한다. 제품 규모의 앱 개발이 압축되지 못한 이유는 백엔드 작업이 압축된 이유의 정확한 거울상이다.

\begin{itemize}
    \item \textbf{피드백 루프가 사람을 거친다.} 화면이 옳게 느껴지는지, 온보딩 흐름이 사람들을 잃는지, 마흔 개 화면에 걸쳐 디자인 언어가 일관되는지, 이런 판정은 종료 코드가 아니라 사용자와 취향에서 온다. 에이전트는 그 어느 때보다 빠르게 후보 UI를 만들어 내지만, 병목은 창업자의 주의력이며 그것은 확장되지 않는다.
    \item \textbf{제품 표면적은 끝이 열려 있다.} 백엔드에는 유한한 계약이 있지만, 소비자 앱은 다듬기 · 플랫폼 업데이트 · 스토어 심사 주기 · 기기별 동작의 긴 꼬리에 대한 무한한 약속이다. 백엔드의 충분함은 통과하는 배포이고, 앱의 충분함은 리텐션 곡선이다.
    \item \textbf{1인의 맥락 전환이야말로 진짜 비용이다.} 생성이 사실상 공짜여도, 유료 고객을 위한 프로덕션 백엔드 운영과 디자인 주도 제품 반복 사이를 오가는 한 사람은 복리로 불어나는 전환 비용을 치른다. 인터럽트 구동이며 검증 가능한 백엔드 작업은 파편화를 견디지만, 창조적 제품 작업은 견디지 못한다.
\end{itemize}

전략적 대응은 비대칭과 싸우기를 그만두는 것이었다. 문AI의 사업 모델은 도구의 모양 그 자체다. 그것은 열린 결말의 앱 개발이 아니라 \emph{백엔드 중심의 AI 시스템}(AI 진단, 레퍼런스 스택 위의 6주 MVP 스프린트, 릴리스 엔지니어 리테이너)을 판다. 내가 온전히 만들어 낼 수 없었던 제품들은 매출 라인이 아니라 역량 증명 사례가 되었다. 도구가 어느 절반을 만들 수 있게 했는지 알고, 오직 그 절반만 짓는 것, 그것이야말로 이 시기가 낳은 핵심 지식이다.

\section{ZUZU: 플랫폼이 된 법인 운영}

\subsection{행정의 벽}

엔지니어링은 문제의 절반일 뿐이었다. 미국 거주 창업자가 운영하는 한국 시장 사업은 역사적으로 전문 인력을 요구해 온 법인 행정 부담에 부딪힌다. 한국 법인의 설립과 유지, 주주명부 관리, 회사의 등기 사항이 바뀔 때마다의 등기 신청, 그리고 한국 은행 · 정부 보조 사업 · 기업 고객이 거래에 앞서 요구하는 서류 이력의 생산이다. 전통적 모델에서 이것은 신청마다 법무사, 원장 사건마다 회계사, 그리고 이상적으로는 이들을 조율할 행정 직원을 뜻한다. 이는 매출 이전 1인 회사가 감당할 수 없는 인건비이자, 열두 시간대 떨어진 창업자가 흡수할 수 없는 조율 부하다.

이 벽은 부수적인 것이 아니다. 한국의 보조 생태계(한국 중소기업에게 AI 프로젝트를 감당 가능하게 만드는 스마트공장 · 바우처 사업)는 등록된 국내 공급기업을 통해 흐른다. 공급기업 등록은 실재하고, 준법하며, 현행인 한국 법인을 요구한다. 준법 법인이 없으면 시장 접근도 없다. 문AI에게 법인 행정은 따라서 간접비가 아니라, 전체 시장 진입 전략(GTM)으로 가는 관문이었다.

\subsection{ZUZU가 흡수한 것}

ZUZU(\href{https://zuzu.network}{zuzu.network})는 이 행정 계층을 소프트웨어로 바꾸는 한국의 법인 운영 플랫폼이다. 설립 워크플로, 주주명부 관리, 등기 신청, 스톡옵션 및 지분 기록, 그리고 이들 각 사건이 요구하는 법적 문서의 생성이다. 창업자의 자리에서 그것을 쓰는 경험은 클라우드 공급자를 쓰는 경험과 매우 흡사하다. 예전에는 채용을 요구하던 기능이 온디맨드로 호출되고, 사용량에 따라 과금되며, 창업자가 그 내부에 정통해질 필요 없이 올바르게 실행된다.

세 가지 성질이 특히 이 국경 횡단 사례에 대해 그것을 충분하게 만들었다.

\begin{itemize}
    \item \textbf{기본이 비동기이며 원격.} 신고와 원장 사건은 서울에서의 대면 회의가 아니라 플랫폼을 통해 준비 · 검토 · 실행된다. 캘리포니아에 거주하는 창업자에게 이것은 성립하는 한국 법인과 불가능한 한국 법인 사이의 차이다.
    \item \textbf{구성상 올바른 서류.} 한국 법인 신고는 형식 오류에 가차 없고, 반려 주기는 주 단위로 측정된다. 요건을 소프트웨어에 인코딩하는 것은 비전문가 창업자가 끊임없이 만들어 낼 부류의 실패를 제거한다.
    \item \textbf{생태계에 대한 가독성.} 법인의 기록이 표준적이고 현행이며 전문적으로 읽히는 형태로 유지되기에, 하류의 상호작용(은행 계좌, 보조금 신청, 공급기업 등록, 훗날의 투자자 실사)이 교정 프로젝트가 아니라 깨끗한 서류 상태에서 출발한다.
\end{itemize}

Claude Code와의 평행은 정확하며, 이것이 이 논문의 논지다. \textbf{ZUZU가 법인 운영 기능에 대해 갖는 관계는, Claude Code가 백엔드 엔지니어링 기능에 대해 갖는 관계와 같다.} 회사가 예전에는 인력으로 채워야 했던 역할 하나를 통째로 흡수하는 플랫폼이다. 미국 쪽에서도 같은 패턴이 범용 부품으로 성립했다. 캘리포니아 단일 구성원 LLC, 상업용 등록 대리인, 그리고 EIN이 며칠 만에 몇백 달러로 준법 청구 수단을 조립했다. 어느 관할권도 나에게 누군가를 채용하라 요구하지 않았다.

\section{조연: 관리형 모델, 결제, 로컬 추론}

ZUZU와 Claude Code는 그러지 않았다면 채용을 요구했을 두 기능, 곧 법인 운영과 백엔드 엔지니어링을 짊어졌다. 그러나 5년치 기록은 그 아래에 놓인, 그야말로 그 위에서 개발하기에 놀라웠던 인프라 계층을 명명하지 않고서는 불완전하다. 세 범주가 명시적 공로를 인정받을 만하다.

\subsection{관리형 모델 API: GCP Vertex AI와 Amazon Bedrock}

문AI가 출시한 모든 AI 기능은 관리형 모델 API 위에서 돈다. GCP 쪽의 Vertex AI(Gemini)와 AWS 쪽의 Amazon Bedrock이다. 이 서비스들이야말로 1인 회사가 애초에 ``AI 시스템''을 팔 수 있는 이유다. 훈련할 모델도, 프로비저닝할 GPU 플릿도, 새벽 3시에 살려 둘 서빙 스택도 없다. 머신러닝 인프라 기능 전체가 토큰당 과금되는 엔드포인트로 도착한다. 특히 Vertex AI는 프로덕션에서 탁월했다. Gemini 모델 계열이 플래그십 카탈로그의 시맨틱 검색 · 채팅 · 설명 생성을 구동하며, 모델 티어(채팅에 \texttt{flash}, 검색 스코어링에 \texttt{flash-lite})가 기능별로 비용을 조율하게 한다. Bedrock은 AWS에서 같은 모양을 제공하며, 이것이 바로 ``하나의 플레이북, 두 클라우드'' 원칙을 실재하게 한다. 두 클라우드 모두 거의 동일한 통합 조작감을 지닌 일급 관리형 모델 서비스를 제공하기에, 플레이북은 고객에게 어느 클라우드든 약속할 수 있다. 그에 못지않게 중요하게도, 토큰당 과금이야말로 애초에 AI 프로젝트를 한국 중소기업 예산 안으로 끌어들인 것이다. 문AI 시장의 경제적 전제는 이 두 서비스가 존재하기에 존재한다.

\subsection{결제: Stripe}

상업적 측면에서 Stripe는 ZUZU가 법인 기능에 한 일을 결제 기능에 했다. 캘리포니아 LLC와 그 EIN 외에는 아무것도 없이, Stripe는 청구, 카드 및 계좌이체 수납, 국경 간 빌링을 제공했다. 회사가 전통적으로 재무 기능을 길러 운영하는 매출채권 기계장치 전체다. 미국 법인에서 청구하면서 두 나라에 걸친 관계로 판매하는 창업자에게, 결제가 API로 도착하고 그 밑에서 준법 · 세금 서식 · 지급 배관이 처리된다는 것은 편의가 아니라 전제 조건이었다. 그것이 1일차 LLC가 실제로 1일차에 매출을 걷을 수 있었던 이유다. Stripe는 클라우드 공급자들과 같은 목록에 속한다. 하나의 사업 기능 전체가 끝내 인력으로 채워지지 않아도 될 만큼 충분히 좋은 인프라다.

\subsection{로컬 추론: Ollama와 MLX 최적화 모델}

모든 모델 호출이 클라우드에 속하는 것은 아니다. Ollama는 로컬 모델 변형(애플 실리콘의 통합 메모리를 활용하는 MLX 최적화 빌드를 포함)을 실행하여 창업자 자신의 기계를 한계 비용 0의 추론 환경으로 바꾸었고, 토큰당 과금이 억제하는 바로 그 작업에 환상적이었다. Vertex AI나 Bedrock에 커밋하기 전 프롬프트와 파이프라인의 빠른 프로토타이핑, 두 나라를 오가며 하는 오프라인 개발, 그리고 첫 질문이 흔히 ``우리 데이터가 건물을 떠나야 하느냐''인 프라이버시 민감 한국 중소기업 잠재 고객을 위한 데모다. 이 패턴은 더 작은 규모에서 이 논문의 나머지를 비춘다. 로컬 추론은 관리형 서비스를 대체하지 않았지만(프로덕션 트래픽은 Vertex AI와 Bedrock에서 돈다), 실험 기능을 흡수하여 비싼 계층을 반복해서 부딪히는 대상이 아니라 의도적으로 도달하는 대상으로 만들었다.

\section{결합: 세 요소가 함께여야만 작동하는 이유}

각 요소는 홀로는 불충분하며, 그렇기에 이 5년의 경험이 자명한 계획이 아니라 하나의 발견처럼 느껴졌다.

\begin{itemize}
    \item \textbf{ZUZU 없는 Claude Code}는 그것을 통해 팔 준법 법인이 없는 배포된 시스템을 낳는다. 특히 한국 시장에서, 돌아가는 법인 없이 돌아가는 소프트웨어는 상업적으로 무력하다.
    \item \textbf{Claude Code 없는 ZUZU}는 제품 없이 완벽하게 관리되는 회사를 낳는다. 창업자 홀로는 듀얼 클라우드 백엔드 포트폴리오를 수용 가능한 속도로 손수 납품할 수 없었을 것이기 때문이다.
    \item \textbf{이중 국적-거주 발판 없는 두 도구}는 사업을 태평양 한쪽에 좌초시킨다. 한국 국적 없이는 한국 법인 · 보조금 접근 · 중소기업 신뢰 구축이 극적으로 어려워지고, 미국 영주권 없이는 창업자가 지금 사는 곳에 합법적으로 거주할 수도, 초기 매출을 미국 LLC에 정박시킬 수도, 납품 역량이 나온 미국 전문가 네트워크를 유지할 수도 없다.
\end{itemize}

결합되면 이들은 다음과 같은 조직도를 가진 회사를 이룬다. 한 명의 인간(판단, 영업, 아키텍처, 취향), 한 개의 에이전틱 엔지니어링 플랫폼(GCP와 AWS 위의 백엔드 납품), 관할권당 한 개의 법인 운영 플랫폼(한국의 ZUZU, 미국의 LLC + 등록 대리인 스택), 그리고 두 여권만큼의 법적 지위(더 정확히는 하나의 여권과 하나의 영주권 카드). 그 구조는 실제 고객을 위한 프로덕션 시스템을 두 클라우드에서 납품하면서도, 인원 수로는 단 한 사람으로 남았다.

\section{교훈}

\begin{enumerate}
    \item \textbf{올바른 기준은 역량이 아니라 충분함이다.} 정작 중요했던 질문은 ``AI가 코딩할 수 있는가?''도 ``소프트웨어가 서류를 낼 수 있는가?''도 아니라, ``그 조합이 추가 채용 0으로 고객의 기준선을 통과하기에 충분한가?''였다. 백엔드 시스템 소프트웨어 더하기 한국 법인 운영에 대해, 2026년에 이르러 그 답은 예가 되었다. 그것은 임계에 관한 진술이며, 사업 모델을 바꾸는 것은 임계다.
    \item \textbf{어느 기능이 인프라로 흡수되었는지 알고, 그 경계를 따라 회사를 빚어라.} 백엔드 엔지니어링과 법인 행정은 선을 넘었고, 앱 개발 · 디자인 · 영업은 넘지 않았다. 문AI의 제품 카탈로그(AI 진단, 백엔드 중심 MVP 스프린트, 릴리스 엔지니어 리테이너)는 정확히 그 경계를 따라 그어졌다. 인프라가 흡수하는 것을 중심으로 제안을 빚는 1인 창업자는 신뢰받을 수 있고, 그 경계가 없는 척하는 창업자는 반쯤 만들어진 앱을 출시한다.
    \item \textbf{검증 가능성이 어느 작업이 압축될지 예측한다.} 백엔드 작업이 압축되고 앱 작업이 그러지 못한 깊은 이유는, 에이전트가 환경 자신이 산출물을 판정하는 곳에서 번성하기 때문이다. 에이전틱 시대에 작은 회사가 무엇을 팔아야 하는지 고를 때, ``성공이 얼마나 기계적으로 검증 가능한가''는 ``그 일이 얼마나 어려운가''보다 나은 길잡이다.
    \item \textbf{이중 관할권 지위는 간접비가 아니라 자산이다.} 여러 해 동안 국민-더하기-영주권 위치는 이중 과세 서류처럼 보였다. 실제로 그것은 해자였다. 한국 국내 공급기업이면서 동시에 미국 거주 운영자일 수 있는 사람은 극히 드물고, 그 위치는 경쟁자가 빠르게 획득할 수 없다.
    \item \textbf{창업자에게 남은 일은 판단이며, 그것은 지켜져야 한다.} 플랫폼이 흡수한 모든 기능은 시간을 돌려주었지만, 그 시간은 오직 그것들이 흡수할 수 없는 기능(아키텍처 결정, 보안 자세, 가격 정책, 고객 신뢰, 그리고 무엇을 짓지 \emph{않을지} 아는 것)으로 흘러 들어갈 때에만 의미가 있었다. 1인 회사는 창업자의 일을 없애지 않는다. 그것을 정련한다.
\end{enumerate}

\section{한계와 남은 질문}

이것은 1인 창업자, 단일 회사의 경험 보고서이며, 일반화에는 주의가 필요하다. 창업자는 구글에서의 3년간(2022--2025) 소프트웨어 엔지니어링을 가지고 문AI에 들어왔다. 여기서 서술한 에이전틱 레버리지는 기존 판단을 증폭할 뿐, 그것을 무(無)에서 부팅하지는 못할 수 있다. 앱 개발 격차는 2026년 시점의 서술이다. 디자인 주도 제품 작업을 위한 에이전틱 도구는 개선되고 있으며, 위 교훈에서 그은 납품 가능성의 경계는 자연법칙이 아니라 움직이는 프런티어로 다뤄져야 한다. 법인 운영 주장은 실제로 겪은 한국과 캘리포니아 제도에 한정된다. 다른 관할권에는 ZUZU에 상응하는 것이 없을 수 있고, 이중 지위의 이점은 정의상 온디맨드로 복제되지 않는다. 끝으로 두 충분함 주장 모두 플랫폼 리스크를 물려받는다. 두 플랫폼 위에 세워진 1인 회사는 인건비를 줄인 만큼이나 의존성을 집중시켰으며, 어느 계층의 가격이나 정책 변화든 회사의 원가 기반으로 곧장 전파된다.

\section{결론}

대한민국과 미국에 걸쳐 2025년에 개발하고 2026년에 출범한 문AI를, 구글에서의 3년간 소프트웨어 엔지니어링을 토대로 홀로 만들면서, 나는 하나의 압축된 사실을 배웠다. 2026년의 1인 창업자는 GCP와 AWS에서 프로덕션 백엔드 시스템 소프트웨어를 납품하고 두 나라에서 준법 법인을 운영할 수 있는데, 그것은 두 범주의 플랫폼(Claude Code로 대표되는 에이전틱 코딩 도구와 ZUZU로 대표되는 법인 운영 플랫폼)이 예전에는 인원 수였던 기능들을 흡수했기 때문이다. 그 아래에서, 똑같이 주목할 만한 기반이 머신러닝 · 재무 · 프로토타이핑 기능을 같은 방식으로 흡수했다. 관리형 모델의 Vertex AI와 Bedrock, 결제의 Stripe, 그리고 실험의 MLX 최적화 로컬 모델을 갖춘 Ollama다. 같은 경험이 그 사실의 경계 또한 가르쳐 주었다. 제품 포트폴리오를 위한 완전한 애플리케이션 개발은 끝내 한 사람의 범위 밖에 남았고, 지속 가능한 수는 그 경계를 가로지르는 것이 아니라 그것을 따라 사업을 설계하는 것이었다. 1인 · 양국 소프트웨어 회사는 더 이상 묘기가 아니다. 그것은 창업자의 법적 지위가 시장들에 걸쳐 있고, 엔지니어링이 에이전트에게 충분히 검증 가능하며, 서류 작업이 API가 되었을 때 벌어지는 일이다.

\end{document}
